\chapter{Talente}
\label{app:talents}

Dieser Anhang enthält alle allgemeinen Talente und Rollentalente, die Helden erlernen können.

\begin{kurzregel}[title=Talente einsetzen]
Nennt ein nicht dauerhaftes Talent weder eine Aktion noch eine Reaktion, wird es im eigenen Zug als \freeAction{} eingesetzt. Eine Reaktion ist eine \freeAction{} außerhalb des eigenen Zuges und darf nur unmittelbar nach ihrem fest definierten Auslöser eingesetzt werden. Alle angegebenen Voraussetzungen und Kosten müssen vollständig erfüllt sein.
\end{kurzregel}

\section{Allgemeine Talente}

\subsection{Magietalente}

\begin{hptable}{L{2.8cm} L{2.7cm} L{2.1cm} Y}
\textbf{Name} & \textbf{Voraussetzung} & \textbf{Einsatz} & \textbf{Effekt} \\
\hline
Magiebegabt & CHA 1+ & Dauerhaft & Du darfst Zauber erlernen und wirken. Dein Manavorrat erhöht sich um 6. \\
Manavorrat & Magiebegabt & Dauerhaft & Dein Manavorrat erhöht sich um 2. \\
Schnelles Zaubern & Magiebegabt & Dauerhaft & Wähle beim Trainieren ein Schlüsselwort. Die Zauberzeit aller Zauber mit diesem Schlüsselwort wird im Kampf um eine Aktionsstufe beschleunigt: \standardAction{} wird zu \moveAction{}, \moveAction{} wird zu \freeAction{}. Eine \freeAction{} kann nicht weiter beschleunigt werden. Außerhalb eines Kampfes wirkst du diese Zauber 30\,\% schneller. \\
Sparsames Zaubern & Magiebegabt & Dauerhaft & Wähle beim Trainieren ein Schlüsselwort. Für alle Zauber mit diesem Schlüsselwort musst du beim Wirken 1 Mana weniger ausgeben, mindestens jedoch 1 Mana. Trainierst du das Talent neu, darfst du ein neues Schlüsselwort wählen. \\
\end{hptable}

\subsection{Körper und Bewegung}

\begin{hptable}{L{2.8cm} L{2.7cm} L{2.1cm} Y}
\textbf{Name} & \textbf{Voraussetzung} & \textbf{Einsatz} & \textbf{Effekt} \\
\hline
Akrobat & GEW 2+ & Dauerhaft & Vorteil auf Proben mit \emph{Turnen} beim Springen, Balancieren oder kontrollierten Fallen. \\
Sprinter & GEW 2+ & \exhaust & Erhalte sofort eine zusätzliche \moveAction{}, die du unmittelbar nutzen darfst. \\
Standfest & KON 2+ & Reaktion + \exhaust & Würdest du den Statuseffekt \emph{Zu Boden} erhalten, darfst du als Reaktion die Talentkarte erschöpfen und verhinderst den Statuseffekt. \\
Schnell auf den Beinen & GEW 1+ & 1 \crystal & Stehe sofort auf, ohne dafür eine \moveAction{} auszugeben. \\
\end{hptable}

\subsection{Wahrnehmung und Natur}

\begin{hptable}{L{2.8cm} L{2.7cm} L{2.1cm} Y}
\textbf{Name} & \textbf{Voraussetzung} & \textbf{Einsatz} & \textbf{Effekt} \\
\hline
Spurenleser & INT 1+ & Dauerhaft & Vorteil auf Proben mit \emph{Natur}, um Spuren zu finden, zu verfolgen oder daraus Informationen abzuleiten. \\
Orientierungssinn & INT 1+ & Dauerhaft & Du kennst unter gewöhnlichen Bedingungen stets die Himmelsrichtungen und kannst einen selbst zurückgelegten Weg wiederfinden. Ist dennoch eine Probe zur Orientierung nötig, hast du Vorteil. \\
Tierversteher & CHA 1+ & Dauerhaft & Vorteil auf Proben mit \emph{Natur}, um Tiere zu beruhigen, ihr Verhalten zu deuten oder mit ihnen umzugehen. Ihre grundsätzliche Stimmung und offensichtlichen Bedürfnisse erkennst du ohne Probe. \\
\end{hptable}

\subsection{Soziales}

\begin{hptable}{L{2.8cm} L{2.7cm} L{2.1cm} Y}
\textbf{Name} & \textbf{Voraussetzung} & \textbf{Einsatz} & \textbf{Effekt} \\
\hline
Menschenkenntnis & KLU 2+ & Dauerhaft & Vorteil auf Proben mit \emph{Herausfinden}, um Absichten, Gefühle oder Lügen einer Person zu erkennen. Eine offensichtliche Stimmung erkennst du ohne Probe. \\
Wortgewandt & CHA 2+ & Dauerhaft & Vorteil auf Proben mit \emph{Überreden}, wenn dein Gegenüber bereit ist, dir zuzuhören und ausreichend Zeit für ein Gespräch besteht. \\
Einschüchternd & MUT 2+ & Dauerhaft & Vorteil auf Proben mit \emph{Überreden}, wenn du eine glaubwürdige Drohung aussprichst oder deine Überlegenheit demonstrierst. \\
Bühnenprofi & CHA 2+ & Dauerhaft & Wähle beim Trainieren eine Form des Auftretens, etwa Musik, Schauspiel, Tanz oder Erzählkunst. Vorteil auf passende Proben mit \emph{Vorführen}. \\
Mitreißend & Bühnenprofi & \exhaust & Nachdem dir eine Probe mit \emph{Vorführen} gelungen ist, darfst du die Talentkarte erschöpfen. Bis zu CHA Personen, die deinen Auftritt miterlebt haben, erhalten Vorteil auf ihre nächste Probe mit einer Fertigkeit der Kategorie \emph{Miteinander} in dieser Szene. \\
\end{hptable}

\subsection{Handwerk, Heilkunde und Wissen}

\begin{hptable}{L{2.8cm} L{2.7cm} L{2.1cm} Y}
\textbf{Name} & \textbf{Voraussetzung} & \textbf{Einsatz} & \textbf{Effekt} \\
\hline
Handwerksmeister & FIN 2+ & Dauerhaft & Wähle beim Trainieren ein Handwerk. Vorteil auf passende Proben mit \emph{Basteln}, um entsprechende Gegenstände herzustellen, umzubauen oder zu reparieren. \\
Improvisateur & Handwerksmeister & \standardAction{} + \exhaust & Stelle aus greifbaren Materialien einen einfachen provisorischen Gegenstand her. Er ermöglicht eine klar bestimmte Verwendung und wird anschließend unbrauchbar. Die Spielleitung entscheidet, was mit den vorhandenen Materialien möglich ist. \\
Ersthelfer & KLU 2+ & \standardAction{} + \exhaust & Ein benachbarter kampfunfähiger Held heilt 1 Lebenspunkt und ist dadurch nicht länger kampfunfähig. \\
Erfahrener Heiler & Ersthelfer & 1 \crystal & Wenn du einer Kreatur Lebenspunkte zurückgibst, darfst du 1 \crystal{} ausgeben. Die Kreatur heilt 1 zusätzlichen Lebenspunkt. \\
Gelehrter & KLU 2+ & Dauerhaft & Wähle beim Trainieren ein Wissensgebiet. Vorteil auf passende Proben mit \emph{Herausfinden}. \\
Geistesblitz & Gelehrter & \exhaust & Nachdem eine Probe mit \emph{Herausfinden} misslungen ist, darfst du die Talentkarte erschöpfen. Du erhältst von der Spielleitung einen wahren, hilfreichen Hinweis. Der Hinweis löst das Problem nicht unmittelbar, zeigt aber einen möglichen nächsten Ansatz. \\
\end{hptable}

\clearpage
\section{Rollentalente}

\subsection{Draufgänger}

\begin{hptable}{L{2.8cm} L{2.7cm} L{2.1cm} Y}
\textbf{Name} & \textbf{Voraussetzung} & \textbf{Einsatz} & \textbf{Effekt} \\
\hline
Perfekter Treffer & -- & 1 \crystal & Nachdem alle Schadenspools eines deiner Angriffe geworfen wurden, aber bevor der Schaden ermittelt wird, darfst du beliebig viele der dabei geworfenen Schadenswürfel auf ihre höchste Seite drehen. \\
Mehrfachangriff & Perfekter Treffer & 1 \crystal & Nachdem alle Sterne ausgegeben und Schadenswürfel verbessert wurden, aber bevor ein Schadenspool verteilt oder geworfen wird, wähle einen zweiten Gegner in Angriffsreichweite und verdopple den vollständigen Schadenspool. Ohne Spalten wird je eine Hälfte des Pools einem der beiden Gegner zugeteilt. Besitzt der Angriff Spalten, dürfen anschließend alle Würfel des verdoppelten Schadenspools nach den Regeln für Spalten verteilt werden. Schild-Symbole und sonstige Abwehrregeln werden für jeden Gegner separat angewendet. \\
Schwachstelle & Perfekter Treffer & 1 \crystal & Nachdem das Ziel eines Angriffs bestimmt wurde, kannst du für diesen Angriff alle Schild-Symbole dieses Gegners ignorieren. \\
Blutrausch & Perfekter Treffer & Dauerhaft & Hast du höchstens die Hälfte deiner Lebenspunkte, werden nach dem Ausgeben aller Sterne und dem normalen Verbessern der Schadenswürfel alle deine Schadenswürfel um eine weitere Stufe verbessert: \orangeDmgDie{} $\rightarrow$ \redDmgDie{}, \redDmgDie{} $\rightarrow$ \blackDmgDie{}, \blackDmgDie{} bleibt \blackDmgDie{}. Erst danach werden Schadenspools durch Mehrfachangriff oder andere Effekte verändert und gegebenenfalls verteilt. \\
Zielen & Perfekter Treffer & \standardAction & Verzichte auf deinen Angriff. Bei deinem nächsten Angriff in diesem Kampf zählt jede leere Seite deiner \blueDie{} und \greenDie{} als ein \star. \\
\end{hptable}

\subsection{Verteidiger}

\begin{hptable}{L{2.8cm} L{2.7cm} L{2.1cm} Y}
\textbf{Name} & \textbf{Voraussetzung} & \textbf{Einsatz} & \textbf{Effekt} \\
\hline
Beschützen & -- & 1 \crystal & Wähle einen benachbarten Gegner. Dieser muss dich mit seiner nächsten Aktion angreifen. Kann er das nicht, verfällt seine nächste Aktion. \\
Bollwerk & Beschützen & Reaktion & Sobald ein benachbarter Gegner versucht, deine Nähe zu verlassen, darfst du als Reaktion vor Abschluss seiner Bewegung einen Gelegenheitsangriff mit einer aktuell verfügbaren Angriffsoption für Schlagen oder Zaubern ausführen. Stammt die Angriffsoption von einem Zauber, musst du dessen Manakosten bezahlen und alle weiteren Voraussetzungen erfüllen; die Zauberzeit wird durch die Reaktion ersetzt. Der erste erzielte \star{} beendet die Bewegung und steht für den Angriff nicht mehr zur Verfügung. Mit allen weiteren \star{} darfst du nach den normalen Regeln Sterne gemäß der gewählten Angriffsoption ausgeben und Schadenswürfel verbessern. Erzielst du keinen \star, setzt der Gegner seine Bewegung fort. \\
Schildwall & Beschützen & Dauerhaft & Benachbarte Verbündete können bei Angriffen gegen sie deine Schild-Symbole anstelle ihrer eigenen verwenden. \\
Verteidigen & Beschützen & \exhaust & Verdopple deine Schild-Symbole bis zum Beginn deines nächsten Zugs. \\
Unbeugsam & Beschützen & Reaktion + \exhaust & Würdest du auf 0 oder weniger Lebenspunkte sinken, darfst du als Reaktion die Talentkarte erschöpfen. Du sinkst stattdessen auf $2 \times \mathrm{KON}$ Lebenspunkte. \\
\end{hptable}

\subsection{Anführer}

\begin{hptable}{L{2.8cm} L{2.7cm} L{2.1cm} Y}
\textbf{Name} & \textbf{Voraussetzung} & \textbf{Einsatz} & \textbf{Effekt} \\
\hline
Heilende Worte & -- & \standardAction & Heile einen benachbarten Helden um dein CHA an Lebenspunkten. Gibst du währenddessen 1 \crystal{} aus, verdoppelt sich die Heilung auf $2 \times \mathrm{CHA}$. \\
Ermutigen & Heilende Worte & Reaktion + 1 \crystal & Nachdem ein Verbündeter seine Probe geworfen hat, darfst du als Reaktion 1 \crystal{} ausgeben. Der Verbündete darf alle \blueDie{} dieser Probe neu würfeln. \\
Inspiration & Heilende Worte & Dauerhaft & Gibst du einen \crystal{} aus, darfst du ihn anschließend an einen Verbündeten deiner Wahl weitergeben. Ein Mondkristall darf insgesamt nur einmal weitergegeben werden; wurde er bereits durch Inspiration oder einen anderen Effekt weitergegeben, kann er nicht erneut weitergegeben werden. \\
Erneuern & Heilende Worte & 3 \crystal{} + \exhaust & Mache eine erschöpfte Talentkarte eines Verbündeten bereit. \\
Aufmuntern & Heilende Worte & \exhaust & Richte einen kampfunfähigen Helden mit vollen Lebenspunkten wieder auf. \\
\end{hptable}

\subsection{Taktiker}

\begin{hptable}{L{2.8cm} L{2.7cm} L{2.1cm} Y}
\textbf{Name} & \textbf{Voraussetzung} & \textbf{Einsatz} & \textbf{Effekt} \\
\hline
Taktieren & -- & 1 \crystal & Sieh dir die nächsten drei KI-Einträge eines Gegners an. Entscheide, zu welchem Eintrag die KI springt. \\
Zu Fall bringen & Taktieren & 1 \crystal & Ein benachbarter Gegner erhält den Kampfeffekt \effectKnockdown{} \emph{Zu Fall bringen} und damit den Statuseffekt \emph{Zu Boden}. \\
Umpositionieren & Taktieren & 1 \crystal & Wähle bis zu zwei dir benachbarte Helden. Sie erhalten jeweils eine kostenlose \moveAction{}, die sie sofort nutzen dürfen. \\
Voraussicht & Taktieren & Dauerhaft & Du darfst bei deinen Proben nach dem Wurf beliebig viele \star{} als \moon{} werten. \\
Meisterplan & Taktieren & \exhaust & Alle Helden führen sofort einen zusätzlichen Zug aus. Anschließend führen die Gegner zwei vollständige Gegnerphasen aus. \\
\end{hptable}

\begin{todo}
Die Funktionsweise von \emph{Taktieren} muss mit dem finalen Gegner- und KI-System abgeglichen werden.
\end{todo}